\documentclass[10pt]{article}
\usepackage[margin=0.9in]{geometry}
\usepackage[T1]{fontenc}
\usepackage{mathptmx}
\usepackage{microtype}
\usepackage{amsmath}
\usepackage{xcolor}
\usepackage{enumitem}
\usepackage[colorlinks=true,urlcolor=blue!50!black,linkcolor=blue!50!black]{hyperref}
\newenvironment{rev}{\par\medskip\noindent\color{black!70}\itshape}{\par}
\newcommand{\resp}{\par\smallskip\noindent\textbf{Response.} }
\newcommand{\chg}{\par\smallskip\noindent\textbf{Changes.} }
\setlist{nosep}
\title{\textbf{Response to reviewers}\\\large Round 1}
\author{Kyler Zook}
\date{October 2026}
\begin{document}
\maketitle

We thank the reviewer for a careful and constructive review. We agree with its central recommendation, to narrow the claims before expanding the experiments, and have restructured the manuscript accordingly.
The main text now covers four things: the response decomposition, independent-screen evaluation, a bounded cascade test, and the limits of reading responses as networks. All simulation branches moved to the Supplement and are stated as results about the specific estimators and regimes tested.

We ran every analysis the reviewer requested that the available data permit. Two of them changed conclusions:
\begin{itemize}
\item the cascade (``domino'') result depends on residualisation, exactly as the reviewer suspected;
\item the context analysis with a noise-matched null shows that context-dependent gains and losses of responses are common, so ``context changes gain, not wiring'' is withdrawn.
\end{itemize}
Section and table numbers below refer to the revised manuscript.

\section*{1. Perturbation effects vs direct edges; the sign-flip analysis}
\begin{rev}Preserved signs of responses do not establish preserved direct edges; conditioning on targets strong in both contexts excludes gained/lost responses; report gained, lost, shared and reversed with a noise-matched null; fixed-sign gated edges should remain a hypothesis.\end{rev}
\resp Agreed on every point. We now test each gene directly for a context difference, $z_\text{diff}=(d_2-d_1)/\sqrt{s_1^2+s_2^2}$. The null is matched to each comparison's sampling noise: both contexts share the precision-weighted mean effect, and each estimate is redrawn with its own standard error.
Significant changes occur at $8.5\times$ the null rate in Norman 2019 and $3.0\times$ in Wessels 2023. Only 1.8\% and 1.1\% of them are reversals; the rest are changes in size (23\% and 5\%), gained responses (26\% and 29\%) and lost responses (25\% and 29\%).
We also show why the dichotomous gained/lost comparison the reviewer described is dominated by power: under the null, ``lost'' calls are as frequent as observed.
\chg Section~2.5, Table~7, Fig.~5. The conclusion is now ``sign preservation among strong shared transcriptional responses, with common context-dependent gains and losses''. The text states that response signs do not imply edge signs. The architectural recommendations (including fixed-sign gated edges) are removed from the main text. The first-version sign-flip table is kept only as superseded (Supplementary Table~S9).

\section*{2. ``Noise-free'' target}
\begin{rev}Independence does not remove noise; the screens differ in timing (day 8 vs day 6); explain which dependencies remain; separate within-screen repeatability.\end{rev}
\resp Agreed. The target is renamed the \emph{independent-screen specific-response target}. Section~2.2 now states:
\begin{itemize}
\item what independence removes: shared sampling noise, so a model cannot be rewarded for fitting the noise realisation it was trained on;
\item what it does not remove: noise in the target itself, and biological differences between the screens (including the timing difference, which we now cite);
\item which dependencies remain: the same cell line, very probably the same guide designs (guide identities are not in the released pseudobulk files, so we could not verify this), and the same processing.
\end{itemize}
Within-screen repeatability can only be estimated from 19 second-guide pairs (residual median $r=0.05$). We report this and state that it cannot be cleanly separated from cross-screen agreement in these data.
\chg Title, abstract, Introduction, Section~2.2, Discussion, Limitations.

\section*{3. Positive control for the domino analysis}
\begin{rev}Demonstrate recovery of a known cascade contribution in matched simulated data; compare raw vs residualised targets; uncertainty and detection limits; otherwise conclude only that this linear reconstruction fails.\end{rev}
\resp Done, and the comparison changed the conclusion. On raw profiles and on generic-removed profiles, the cascade prediction clearly beats random mediators: $+0.137$ [0.119, 0.157] and $+0.133$ [0.112, 0.153]. On the specific residual it does not: $-0.004$ [$-0.012$, $0.004$]. Cascade-like structure therefore lies along the shared axes that residualisation removes, as the reviewer anticipated. We note that this structure is equally consistent with shared module membership.

We built a semi-synthetic positive control matched to the real mediator set, the real weights $a_{jk}/d_k$, the real mediator profile shapes and each profile's real cross-screen reliability. It varies the cascade share $\alpha$ and mediator coverage. The identical pipeline detected $\alpha=0.05$ at every coverage level tested (down to 25\%). The real-data upper bound lies below the $\alpha=0.05$ lower bound even at 25\% coverage.

The bounded conclusion: within an additive, linear model, a cascade share $\ge5\%$ of the reproducible specific response is excluded if at least a quarter of each target's true mediators were knocked down. Cascades through never-perturbed genes, and nonlinear relays, are not excluded. Path double-counting is acknowledged.
The reliability-normalised ``domino share'' turned out to be badly attenuated (0.11 at $\alpha=1$), so it is no longer used as an effect size.
\chg Section~2.3, Tables~4--5, Fig.~3; abstract.

\section*{4. Reliability calculations}
\begin{rev}Cross-screen correlation is not a ceiling; under independent noise replicate correlation is $\sqrt{R_1R_2}$ and the maximum is $\sqrt{R_2}$; ``a third of reliability'' compared a mean with a median.\end{rev}
\resp Correct, and the correction makes our predictors look worse. Section~2.2 now derives $\operatorname{corr}(x_1,x_2)=\sqrt{R_1R_2}$ and $\max\operatorname{corr}(f,x_2)=\sqrt{R_2}$. It states that $\sqrt{r_{12}}$ is the ceiling only under equal reliability, and is conservative if the screens' true responses differ.
The ceiling is estimated for the same targets and the same gene panel as the scores. The best predictor reaches about 11\% of it on the specific target, not ``a third''.
\chg Section~2.2, Table~3, Fig.~2.

\section*{5. Decomposition: interpretation and validation}
\begin{rev}A mean direction, PCs and residual are a statistical decomposition; the noise model is transferred from controls; Fig.~3a mixed medians and pooled fractions; report sensitivity, validate noise calibration, show PC stability/enrichment, distinguish pooled vs median.\end{rev}
\resp Agreed throughout.
\begin{itemize}
\item \textbf{Labels.} The components are now explicitly statistical: ``generic'', ``shared axes'' (no longer ``pathways'') and ``specific''.
\item \textbf{Pooled vs median.} Figure~1a shows a consistently pooled decomposition; the median per-knockdown noise share is reported separately (pooled noise 47\% of energy, median 74\%).
\item \textbf{Sensitivity.} We report sensitivity to $K$ (0--50), the gene panel (all, 2{,}000 most variable, expressed genes), cell count and knockdown strength (Table~1). The residual's share of signal ranges from 0.18 to 0.78 depending on $K$, and we say so plainly.
\item \textbf{Noise calibration.} Split halves of perturbed cells are impossible with the released pseudobulk files. We instead validated the noise model across screens: the gw$-$es difference for the same knockdown is $1.30\times$ the predicted noise energy, against $1.18\times$ for random non-targeting pairs (Table~2). The model is approximately right but probably underestimates noise in perturbed profiles.
\item \textbf{Axis stability.} The generic direction and two of three axes reproduce when fitted separately per screen (cosines 0.91, 0.97, 0.94); the third does not (0.01).
\item \textbf{Biological content.} Top-loading genes are ribosomal (axis 1), mitochondrially encoded (axis 2) and erythroid (axis 3, non-reproducible). We did not perform formal enrichment and say so.
\end{itemize}
\chg Section~2.1, Tables~1--2, Fig.~1.

\section*{6. Benchmark novelty and protocol}
\begin{rev}Specify the contribution relative to Systema; full fold-by-fold protocol (gene selection, projection, $k$, feature availability); compare with mean-centring alone; add a magnitude-sensitive metric.\end{rev}
\resp
\begin{itemize}
\item \textbf{Relation to Systema.} The Introduction and Discussion now credit Systema for establishing the shared-variation problem. Our additions are limited to independent-screen scoring, the reliability derivation and the cascade test.
\item \textbf{Leak fixed.} Reviewing the protocol revealed one: the 1{,}000-gene panel had been chosen on all knockdowns, including test ones. It is now chosen per fold on training knockdowns only, and all results were rerun.
\item \textbf{Protocol.} Section~2.2 gives the complete procedure: per-fold refitting of the generic direction and axes ($K$ 2--4 by permutation null), projection of the independent-screen target with the training-fold components, the similarity feature (computed only from training knockdowns, so available at prediction time), and $k\in\{10,20,50,100\}$ reported in full with no tuning.
\item \textbf{New comparisons.} We added a mean-centred independent-screen target (Systema-like centring) and the magnitude-sensitive $R^2_0$. No predictor beats predicting zero in magnitude on independent-screen targets.
\item \textbf{Discrimination.} We agree that the mean profile scoring zero is expected by construction, and now say so. The target does discriminate among substantive predictors: $k=10$ gives 0.048 [0.040, 0.056] and $k=100$ gives 0.071 [0.062, 0.080].
\end{itemize}
\chg Section~2.2, Table~3, Fig.~2.

\section*{7. Simulation-based recommendations exceed tested conditions}
\begin{rev}Exact regulator masks are privileged; the mask + site prior result conflicts with ``do not weight edges by motif''; propagation could not test deep effects; failure of a lag-feature detector does not generalise; noise thresholds need calibration.\end{rev}
\resp All three examples are addressed, and the architectural-recommendation table is removed.
\begin{itemize}
\item \textbf{Masks.} New imperfect-mask experiments (missing 10--50\% of regulators; adding 50--200\% false regulators; combinations; 10 seeds; Supplementary Table~S6, Fig.~S3) show that each missing regulator costs far more than each false one. Dropping 25\% of regulators already falls to the no-prior level, and a mask with both kinds of error is worse than no prior. A hard mask is worthwhile only with high-recall annotation. We corrected the contradiction: the text now says ``mask first; a site prior adds a little on top''.
\item \textbf{Propagation.} New true-vs-estimated operator experiments cover every regulator, stratified by true downstream depth (Supplementary Table~S7, Fig.~S4). With the true operator, deeper propagation helps where deep cascades exist. With the estimated operator, one step is best and deeper steps fail. The limitation is operator accuracy, not depth.
\item \textbf{Memory.} Conclusions are restricted to the lag-feature ridge detector. Noise levels are described as relative to simulated fluctuations, and the word ``realistic'' is removed.
\end{itemize}
\chg Supplement S5, S7, S8; main-text Limitations.

\section*{8. Binding analysis scope}
\begin{rev}Promoter-specific scope; per-TF uncertainty; characterise the hub baseline; separate failure of occupancy from failure of binding to predict responses.\end{rev}
\resp Section~2.4 states the 1-kb promoter scope up front and reports bootstrap CIs over TFs and paired tests for every score.
The hub score correlates only weakly with ATAC (0.33) and expression (0.23), and keeps its AUROC within accessibility or expression quintiles (0.753, 0.758). It therefore reflects broadly co-occupied promoters or assay structure (for example highly ChIP-able regions), not accessibility or activity.
Separating the two failures: measured ChIP binding predicts responders no better than chance (0.496 [0.473, 0.511]), also within expression quintiles. The failure is therefore not one of the occupancy approximation.
\chg Section~2.4, Table~6, Fig.~4; abstract.

\section*{9. Central argument, methods, reproducibility, baselines}
\resp
\begin{itemize}
\item \textbf{Argument.} The main text is re-centred as the reviewer suggested; the experiment IDs and the narrative of successive hypotheses are gone.
\item \textbf{Methods.} Methods are expanded in the main text (Replogle processing, decomposition, benchmark, cascade test, occupancy calibration, ChIP and hub definitions, context statistics, uncertainty) and in the Supplement (latent ODE architecture and training, the 80-gene simulator, BOND-V, metric definitions).
\item \textbf{Reproducibility.} We provide a versioned repository (tag \texttt{v1.1-revision1}) with \texttt{requirements.txt} (pinned versions), \texttt{data/README.md} (identifiers for every input, including ENCODE accessions), and \texttt{run\_all.sh} (every command, in order, that regenerates each result, table, figure and PDF).
\item \textbf{Baselines.} Re-implemented baselines are described as ``diagnostic re-implementations of the core idea'', and no comparative claims about the original methods remain.
\end{itemize}

\section*{Concrete corrections}
\begin{itemize}
\item \textbf{Table 8 broken:} an unescaped \texttt{\%} in the row label truncated the row. Fixed (now Supplementary Table~S5).
\item \textbf{Table 9, $\infty$ column:} (now Supplementary Table~S8) relabelled $(I-J)^{-1}$, with a footnote that it equals the infinite sum only when $\rho(J)<1$.
\item \textbf{Table 3, zero predictor:} correlation is now reported as undefined (n/a) for constant predictors; $R^2_0$ is reported for all.
\item \textbf{T1:} the two numbers were different comparisons (all targets: 0.959 vs 0.820; weak-cluster targets: 0.968 vs 0.647). Both are now labelled (Supplement S1).
\item \textbf{T3:} the improvement is $\Delta R^2=0.011$ (0.006$\to$0.017); 0.017 is the $R^2$ itself. Corrected.
\item \textbf{Metric definitions:} the candidate-edge universe, AUPRC normalisation and top-$E$ are defined at the start of the Supplement.
\end{itemize}

\end{document}
